\chapterlede{This chapter summarizes what was achieved and what the next steps should be.}
\chapter{Conclusions and future work}
\label{ch:conclusions}
The internship turned a supplied clustering script and data set into an inspectable, tested analysis. The audit identified problems in parsing, residual-cluster handling, FES interpolation, replica pooling and weighted summaries, and the modular reimplementation corrects them with 38 passing tests while preserving the original sources. All three parts of the final task were carried out; \cref{tab:objectives} records the status of each.

The corrected workflow runs the published density-peak clustering on all \num{300003} frames with the published parameters. It finds \DPCFamilies{} families that differ mainly in secondary structure and are shared unevenly between replicas, and the parameter study shows that the partition depends strongly on the seed, the parameters and the data used for discovery. Under an assumed \kJmol{} FES unit, reweighting concentrates \WTopThree\% of the equilibrium weight on three families and almost removes the helix-rich ones, a pattern that would also hold if the file were in \kcalmol{}. The independent Gaussian-mixture experiment shows that the strong $R_g$--SASA correlation does not settle the feature choice: removing SASA largely preserves the three-component partition (ARI 0.862--0.868) but not the others (median ARI 0.457--0.690).

The next steps, in order of priority, are:
\begin{protocolsteps}
\item confirm with the authors the FES unit and how the replicas were combined; the populations of \cref{sec:weighted} can then be recomputed immediately from the saved labels;
\item repeat the seed study with the 20-seed profile that the full configuration already defines, and compare the families with those reported in the primary article;
\item repeat the feature experiment on disjoint time blocks, to measure the sampling sensitivity that the fixed fitting sample deliberately excludes \autocite{grossfield2019};
\item contrast final-bias weights with time-dependent \code{metad.rbias} weights once the metadynamics settings are known \autocite{tiwary2015,bonomi2009}, and compare with DPA \autocite{derrico2021} and with an experimentally constrained Aβ40 ensemble \autocite{borthakur2025};
\item with atomic trajectories and a frame map, extract representative structures and test coordinate-based methods such as ShapeGMM \autocite{michaud2011,sasmal2024}.
\end{protocolsteps}
